\documentclass[12pt,a4paper]{article}

% --- Pacotes Básicos ---
\usepackage[utf8]{inputenc}
\usepackage[portuguese]{babel}
\usepackage[T1]{fontenc}
\usepackage{amsmath}
\usepackage{amsfonts}
\usepackage{amssymb}
\usepackage{geometry}
\usepackage{indentfirst}
\usepackage{listings}
\usepackage{xcolor}

% --- Configuração de Margens ---
\geometry{left=3cm, top=3cm, right=2cm, bottom=2cm}

% --- Configuração para Código C ---
\definecolor{mGreen}{rgb}{0,0.6,0}
\definecolor{mGray}{rgb}{0.5,0.5,0.5}
\definecolor{mPurple}{rgb}{0.58,0,0.82}
\lstset{
    language=C,
    basicstyle=\ttfamily\small,
    commentstyle=\color{mGreen},
    keywordstyle=\color{blue},
    numberstyle=\tiny\color{mGray},
    stringstyle=\color{mPurple},
    breaklines=true,
    showstringspaces=false,
    numbers=left,
    frame=single
}

% --- Título e Autores ---
\title{\textbf{Relatório Primeiro Trabalho Prático da disciplina de Otimização(CI-1238)}}
\author{
    Eduardo Munaretto \\
    Viktor Hugo \\
    João Pedro de Oliveira Lazari (20242306) \\
    \small{UFPR}
}
\date{Abril de 2026}

\begin{document}

\maketitle

\begin{abstract}
 O objetivo do projeto foi determinar a combinação ideal de comprimidos que satisfaça as  necessidades do paciente e ao mesmo tempo respeite os limites de consumo para não prejudicar a saúde, tudo isso com o menor custo possível. A solução envolve modelagem por Programação Linear e uma implementação em C para geração de entradas adequadas para o algoritmo Simplex \texttt{}.
\end{abstract}

\section{Introdução}
Na indústria farmacêutica, é comum que princípios ativos sejam administrados via coquetéis que contêm não apenas os componentes necessários, mas também componentes como lactose, glúten ou álcool. Para pacientes com restrições alimentares , o consumo desses comprimidos deve ser feito com cuidado.

O problema proposto exige que sejam atendidas quantidades mínimas diárias de $c$ componentes a partir de $p$ tipos de comprimidos disponíveis, cada um com um custo específico. Adicionalmente, deve-se respeitar um teto máximo para $k$ componentes limitados.

\section{Modelagem Matemática}

Modelamos o problema por um PL (Programação Linear)

\subsection{Variáveis de Decisão}
Seja $x_i$ a quantidade de comprimidos do tipo $i$ a serem consumidos por dia, para $i = 1, \dots, p$. 

\subsection{Função Objetivo}
Minimizar o custo total da dieta medicamentosa:
\[ \min Z = \sum_{i=1}^{p} v_i x_i \]
$v_i$ é o preço do comprimido $i$.

\subsection{Restrições}
\begin{enumerate}
    \item \textbf{Necessidades Mínimas:} Para cada componente $j \in \{1, \dots, c\}$:
    \[ \sum_{i=1}^{p} r_{i,j} x_i \ge q_j \]
    Onde $r_{i,j}$ é a quantidade do componente $j$ no comprimido $i$, e $q_j$ é a exigência diária.
    
    \item \textbf{Limites Máximos:} Para cada componente limitado $f_t$ com limite $l_t$:
    \[ \sum_{i=1}^{p} r_{i,f_t} x_i \le l_t \]
    
    \item \textbf{Não-negatividade e Integralidade:}
    \[ x_i \ge 0 \quad \text{e} \quad x_i \in \mathbb{Z} \]
\end{enumerate}

\subsection{Justificativa da Modelagem}
 Ao definir as variáveis como inteiras, garantimos que a solução seja aplicável no mundo real (não se pode tomar frações de comprimidos).

\section{Implementação}
A implementação foi realizada em linguagem C.

\subsection{Estrutura do Código}
O programa foi dividido em módulos:
\begin{itemize}
    \item \texttt{printFuncaoObjetivo}: Varre os dados de entrada para montar a linha de custo.
    \item \texttt{printComponenteDiario}: Gera as restrições do tipo $\ge$ para os princípios ativos.
    \item \texttt{printComponenteLimite}: Gera as restrições do tipo $\le$ para os componentes com restrição de consumo.
    \item \texttt{printDeclaracoes}: Define as variáveis como inteiras para o solver.
\end{itemize}


\section{Conclusão}


\end{document}
